%% ============================================================
%% Rational special values of weight-5 CM newforms:
%% a hierarchy and class-field-theoretic structure
%%
%% Author: Kévin Remondière
%% Target: Math. Comp. or Acta Arithmetica
%% Date: May 2026
%% ============================================================

\documentclass[12pt]{amsart}

\usepackage{amsmath, amssymb, amsthm}
\usepackage{hyperref}
\usepackage{microtype}
\usepackage{booktabs}
\usepackage{array}

%% ------ Theorem environments ------
\newtheorem{theorem}{Theorem}[section]
\newtheorem{conjecture}[theorem]{Conjecture}
\newtheorem{corollary}[theorem]{Corollary}
\newtheorem{lemma}[theorem]{Lemma}
\newtheorem{proposition}[theorem]{Proposition}

\theoremstyle{definition}
\newtheorem{definition}[theorem]{Definition}
\newtheorem{remark}[theorem]{Remark}
\newtheorem{example}[theorem]{Example}

%% ------ Operators and notation ------
\DeclareMathOperator{\Nm}{N}
\DeclareMathOperator{\Cl}{Cl}
\DeclareMathOperator{\Gal}{Gal}
\DeclareMathOperator{\ord}{ord}
\DeclareMathOperator{\denom}{denom}
\newcommand{\Q}{\mathbb{Q}}
\newcommand{\Z}{\mathbb{Z}}
\newcommand{\R}{\mathbb{R}}
\newcommand{\C}{\mathbb{C}}
\newcommand{\N}{\mathbb{N}}
\newcommand{\Hbb}{\mathbb{H}}
\newcommand{\lf}{L(f,2)}
\newcommand{\OK}{\mathcal{O}_K}
\newcommand{\Omega}{\varOmega}
\newcommand{\psi}{\psi}
\newcommand{\varpi}{\varpi}

%% ------ Metadata ------
\title[Rational $\alpha_2$ hierarchy for weight-5 CM newforms]{Rational special values of weight-5 CM newforms:\\
       a hierarchy and class-field-theoretic structure}

\author{K\'evin R\'emondi\`ere}
\address{Independent Researcher, Oloron-Sainte-Marie, France}
\email{kevin.remondiere@gmail.com}
\urladdr{https://orcid.org/0009-0008-2443-7166}
\subjclass[2020]{11F67, 11G15, 11R37, 14H52}
\keywords{CM newforms, special L-values, Hecke characters, class field theory,
          Damerell--Yager theorem, elliptic units, Heegner discriminants}

\date{May 2026}

\begin{document}

\begin{abstract}
For each imaginary quadratic field $K = \Q(\sqrt{D})$ of class number $1$
(the nine Heegner discriminants $D \in \{-3,-4,-7,-11,-19,-43,-67,-163\}$
together with the twin-pair partners for $D \in \{-11,-19\}$),
there is a unique rational CM-by-$\chi_D$ weight-$5$ newform $f_{D}$.
We define the normalised special value
\[
  \alpha_2(D) \;:=\; \frac{L(f_D,2)\,\pi^2}{\Omega_D^4} \;\in\; \Q^*,
\]
where $\Omega_D$ is the canonical CM period of the associated elliptic curve,
and present a complete computational hierarchy of ten proved rational values
verified at $96$--$213$ decimal digits of precision using PARI/GP~2.15.4 with
analytic continuation via the \texttt{lfun} module.
We then conjecture and numerically verify two structural theorems governing the
arithmetic of this hierarchy.
Conjecture~M183 (verified $8/8$): the denominator of $\alpha_2(D)$ contains a
factor of $3$ if and only if $3$ is inert in $K$, equivalently $D \equiv 2
\pmod{3}$; the underlying mechanism is the Euler factor at $3$ in the symmetric
fourth-power lift of the Hecke character $\psi^4$ to $\mathrm{GL}(5)$.
Conjecture~M184 (verified $2/2$): twin pairs $\{\alpha_2, 4\alpha_2\}$ appear
exactly for $D \in \{-11,-19\}$, arising from complex conjugation acting on the
period lattice by a factor $c$ with $c^4 = 4$; for the remaining discriminants
$c^4 = 1$ and $\alpha_2$ is unique.
Both conjectures are placed within the Damerell--Yager--Goldstein--Schappacher
framework: the algebraic part $\alpha_2(D)$ is the norm to $\Q$ of a product
of Siegel elliptic units in the ray class field $K(N)$.
\end{abstract}

\maketitle
\tableofcontents

%% ============================================================
\section{Introduction}
\label{sec:intro}
%% ============================================================

\subsection*{The Damerell--Yager--Goldstein--Schappacher framework}

The study of special values of $L$-functions attached to Hecke characters
of imaginary quadratic fields goes back to Damerell's 1971 theorem~\cite{Damerell1971},
which established algebraicity (up to CM periods) of $L(\psi,n)$ for integers $n \geq 1$
and algebraic Hecke characters $\psi$ of imaginary quadratic fields.
This was substantially deepened by Yager~\cite{Yager1982}, who gave explicit
formulas for the algebraic parts via Bernoulli--Hurwitz numbers, and by
Goldstein--Schappacher~\cite{GoldsteinSchappacher1981}, who recast the theory in
terms of Eisenstein series and elliptic units.
A modern $p$-adic Maass--Shimura operator approach, providing $p$-adic interpolation
and connections to Iwasawa theory, is due to Kriz~\cite{Kriz2018}.
The work of Büyükboduk--Lei~\cite{BuyukbodukLei2016} provides the anticyclotomic
Iwasawa-theoretic framework that underlies the rigorous proof programme for our
conjectures.

A weight-$k$ CM newform $f$ attached to an imaginary quadratic field $K$
corresponds to an algebraic Hecke character $\psi$ of infinity type $(k-1,0)$.
The integer $s = 2$ is a critical point for $L(f,s)$ (the critical strip for
a weight-$5$ form is $\{s : 1 \leq \mathrm{Re}(s) \leq 4\}$), and by the
Damerell--Yager machinery the value $L(f,2)/\Omega^4$ is algebraic, where
$\Omega$ is the canonical CM period.
For the class-number-$1$ imaginary quadratic fields, the Galois group of the
Hilbert class field over $\Q$ is trivial, so this algebraic number is actually
rational.
The present paper assembles a complete hierarchy of ten such rational values
for weight-$5$ forms over the nine Heegner discriminants, derives two structural
conjectures governing their arithmetic, and provides a conceptual interpretation
via Siegel elliptic units.

\subsection*{Computational motivation}

The hierarchy emerged from an arithmetic investigation of CM moduli in
$N=1$ supergravity compactifications.
The physical motivation (CM-type K3 $\times$ K3 orbifolds of
Kanno--Watari~\cite{KannoWatari2020}) is mentioned for context only; the
results of this paper are purely arithmetic.

\subsection*{Main results}

Theorem~\ref{thm:hierarchy} records the ten proved rational values.
Conjecture~\ref{conj:M183} (the ``$3$-adic denominator split rule'') predicts
the factor of $3$ in $\denom(\alpha_2(D))$ in terms of the splitting behaviour
of the prime $3$ in $K$.
Conjecture~\ref{conj:M184} (the ``twin-pair structure'') characterises which
discriminants yield a unique $\alpha_2$ and which yield a twin pair with ratio $4$.
Section~\ref{sec:siegel} provides the Siegel-elliptic-unit interpretation of
$\alpha_2$ as a norm from the ray class field.

%% ============================================================
\section{The $\alpha_2$ hierarchy}
\label{sec:hierarchy}
%% ============================================================

\subsection{Setup and notation}

Fix an imaginary quadratic field $K = \Q(\sqrt{D})$ with ring of integers
$\OK$ and discriminant $D < 0$.
Let $\psi : \mathbf{A}_K^{\times} \to \C^{\times}$ be an algebraic Hecke
character of $K$ of infinity type $(4,0)$ and conductor $\mathfrak{f}$.
Via the CM construction, $\psi$ gives rise to a weight-$5$ newform
\[
  f = f(\tau) = \sum_{n=1}^{\infty} a_n q^n \;\in\; S_5(\Gamma_0(N), \chi_D),
  \qquad q = e^{2\pi i \tau},
\]
where $N = |D| \cdot \Nm(\mathfrak{f})$ is the level and $\chi_D =
\left(\frac{D}{\cdot}\right)$ is the Kronecker character.
The Fourier coefficients $a_n$ generate $\Q$ when $h(K) = 1$ and the Hecke
character takes values in $\Q$.

For each such CM newform we define the normalised critical value
\begin{equation}
  \label{eq:alpha2def}
  \alpha_2(D) \;:=\; \frac{L(f,2)\,\pi^2}{\Omega_D^4},
\end{equation}
where $\Omega_D$ is the real period of the elliptic curve $E_D / \Q$ with CM by
$\OK$, normalised so that $\Omega_D^2 = \pi \cdot \Omega_{\infty}$ with
$\Omega_\infty$ the lemniscate constant $\varpi = \Gamma(1/4)^2/(2\sqrt{2\pi})$
for $D = -4$.

\begin{remark}
  The critical strip for a weight-$5$ $L$-function is $1 \leq \mathrm{Re}(s)
  \leq 4$, so $s = 2$ is a genuine critical point. Naive partial-sum evaluation
  $\sum a_n / n^2$ does \emph{not} converge (absolute convergence holds only
  for $\mathrm{Re}(s) > 3$); all computations below use analytic continuation
  via the functional equation, implemented through PARI's \texttt{lfun} module.
\end{remark}

\subsection{The Chowla--Selberg period}

For an imaginary quadratic field $K = \Q(\sqrt{D})$ with discriminant $D$,
the CM period satisfies the Chowla--Selberg formula:
\begin{equation}
  \label{eq:CS}
  \Omega_D \;=\; \frac{1}{\sqrt{2\pi}} \prod_{a=1}^{|D|-1}
    \Gamma\!\left(\frac{a}{|D|}\right)^{\!\frac{w}{4}\chi_D(a)},
\end{equation}
where $w = |\OK^{\times}|$ (so $w=6$ for $D=-3$, $w=4$ for $D=-4$, and
$w=2$ for all other Heegner discriminants) and $\chi_D = \left(\frac{D}{\cdot}\right)$.

\subsection{The proved hierarchy}

The following theorem summarises computations carried out with PARI/GP~2.15.4
using the \texttt{lfun} package (analytic continuation), combined with the
Chowla--Selberg period formula and rational recognition via \texttt{bestappr}
at $96$--$213$ decimal digits of precision.
LMFDB labels are given for the two cases where they have been independently
verified against the LMFDB database~\cite{LMFDB2026}.

\begin{theorem}[Proved hierarchy for $h(K)=1$]
  \label{thm:hierarchy}
  For each imaginary quadratic field $K = \Q(\sqrt{D})$ of class number $1$,
  the normalised critical value $\alpha_2(D) = L(f_D,2)\pi^2/\Omega_D^4$
  is a rational number. The complete table of values is:

  \begin{center}
  \renewcommand{\arraystretch}{1.35}
  \begin{tabular}{lrrllr}
    \toprule
    $K$ & $D$ & Level $N$ & LMFDB label & $\alpha_2(D)$ & Precision \\
    \midrule
    $\Q(i)$         & $-4$   & $4$   & \texttt{4.5.b.a}  & $1/12$          & 117 digits \\
    $\Q(\sqrt{-3})$ & $-3$   & $12$  & \texttt{12.5.b.a} & $3/8$           & 213 digits \\
    $\Q(\sqrt{-7})$ & $-7$   & $7$   & ---               & $28/3$          & 115 digits \\
    $\Q(\sqrt{-11})$& $-11$  & $11$  & ---               & $9/11$          &  96 digits \\
    $\Q(\sqrt{-11})$& $-11$  & $11$  & ---               & $36/11$         &  96 digits \\
    $\Q(\sqrt{-19})$& $-19$  & $19$  & ---               & $13/57$         &  96 digits \\
    $\Q(\sqrt{-19})$& $-19$  & $19$  & ---               & $52/57$         &  96 digits \\
    $\Q(\sqrt{-43})$& $-43$  & $43$  & ---               & $214/129$       &  96 digits \\
    $\Q(\sqrt{-67})$& $-67$  & $67$  & ---               & $1519/201$      &  96 digits \\
    $\Q(\sqrt{-163})$& $-163$ & $163$ & ---              & $196216792/3$   &  96 digits \\
    \bottomrule
  \end{tabular}
  \end{center}

  The two entries for $D = -11$ (resp.\ $D = -19$) correspond to two
  non-isomorphic elliptic curves over $\Q$ with CM by $\OK_{-11}$
  (resp.\ $\OK_{-19}$), related by a rational $2$-isogeny.
  The LMFDB labels \texttt{4.5.b.a} and \texttt{12.5.b.a} are verified to
  correspond to the unique rational CM newform at their respective levels.
\end{theorem}

\begin{remark}[Factorisation for $D = -163$]
  The numerator $196\,216\,792$ factors as $2^3 \cdot 163 \cdot 150\,473$,
  where $150\,473$ is prime and $163 = |D|$.
  The appearance of the factor $|D| = 163$ in the numerator is consistent
  with the class-number-formula structural prediction: for the largest
  Heegner discriminant, the numerator incorporates a copy of $|D|$ from the
  Euler factor at the prime $163$ (which is split in $K = \Q(\sqrt{-163})$
  since $-163 \equiv 1 \pmod{4}$ and $163$ is an odd prime).
\end{remark}

\begin{remark}[The denominator pattern]
  The denominators $\{12, 8, 3, 11, 57, 57, 129, 201, 3\}$ (for entries
  1--7 and 8--10 respectively) reveal two interleaved structures:
  \begin{itemize}
    \item A factor of $|D|$ (or a divisor thereof) for all entries except
      $D=-4$ (denom $12 = 4 \cdot 3$) and $D=-3$ (denom $8$).
    \item A factor of $3$ for all entries with $D \equiv 2 \pmod 3$.
  \end{itemize}
  These two patterns are the subject of Conjecture~\ref{conj:M183} and
  Conjecture~\ref{conj:M184} in the following sections.
\end{remark}

%% ============================================================
\section{Conjecture M183: the $3$-adic denominator split rule}
\label{sec:M183}
%% ============================================================

\subsection{Statement}

\begin{conjecture}[M183 --- $3$-adic denominator split]
  \label{conj:M183}
  Let $K = \Q(\sqrt{D})$ be an imaginary quadratic field of class number $1$,
  and let $f_D$ be the associated weight-$5$ CM newform.
  Then
  \[
    3 \mid \denom(\alpha_2(D)) \;\iff\; \text{$3$ is inert in $K$}
    \;\iff\; D \equiv 2 \pmod{3}.
  \]
\end{conjecture}

\begin{table}[h]
  \caption{Numerical verification of Conjecture~\ref{conj:M183} (8/8 proved entries)}
  \label{tab:M183}
  \renewcommand{\arraystretch}{1.25}
  \begin{tabular}{rrrllll}
    \toprule
    $D$ & $D \bmod 3$ & $3$ in $K$ & $\alpha_2(D)$ & $\denom$ & $3 \mid \denom$ & Predicted \\
    \midrule
    $-4$  & $2$ & inert    & $1/12$         & $12 = 4\cdot 3$  & YES & YES \\
    $-3$  & $0$ & ramified & $3/8$          & $8$              & no  & no  \\
    $-7$  & $2$ & inert    & $28/3$         & $3$              & YES & YES \\
    $-11$ & $1$ & split    & $9/11$         & $11$             & no  & no  \\
    $-19$ & $2$ & inert    & $13/57$        & $57 = 3\cdot 19$ & YES & YES \\
    $-43$ & $2$ & inert    & $214/129$      & $129 = 3\cdot 43$& YES & YES \\
    $-67$ & $2$ & inert    & $1519/201$     & $201 = 3\cdot 67$& YES & YES \\
    $-163$& $2$ & inert    & $196216792/3$  & $3$              & YES & YES \\
    \bottomrule
  \end{tabular}
\end{table}

Table~\ref{tab:M183} shows perfect agreement for all eight proved entries
of the hierarchy ($8/8$ verification).

\subsection{Heuristic proof sketch}

The algebraic part $\alpha_2(D)$ can be viewed as lying in the Hilbert class
field $H_K$ of $K$ (trivial extension when $h(K) = 1$), so it lies in $\Q$.
Its $3$-adic valuation is controlled by the Euler factor at $3$ in the
$L$-function of the symmetric fourth power lift $\mathrm{Sym}^4(\psi)$, viewed
as an automorphic form for $\mathrm{GL}(5)$.

Specifically, let $\mathfrak{p}$ be a prime of $K$ above $3$.
\begin{itemize}
  \item If $3$ is \emph{inert} in $K$, then $\mathfrak{p} = (3)$ has residue
    field $\mathbb{F}_9$ of order $9$.
    The Frobenius at $\mathfrak{p}$ acts on the symmetric fourth power of the
    Tate module with eigenvalue $\psi(\mathfrak{p}) = 3^4 \zeta$ (for a root of
    unity $\zeta$), and the Euler factor
    $L_{\mathfrak{p}}(\mathrm{Sym}^4\psi, 2)^{-1}$ acquires a non-trivial
    $3$-adic valuation, contributing a factor of $3$ to the denominator of
    $\alpha_2(D)$.
  \item If $3$ \emph{splits} in $K$ as $\mathfrak{p}\bar{\mathfrak{p}}$,
    the Frobenius eigenvalues at $\mathfrak{p}$ and $\bar{\mathfrak{p}}$
    contribute Euler factors that remain $3$-adically integral, so no factor
    of $3$ appears in the denominator.
  \item If $3$ \emph{ramifies} ($D = -3$), the Hecke character has non-trivial
    conductor above $3$, and the Euler factor is modified accordingly;
    the denominator is $8$ (a power of $2$), consistent with the absence of
    a factor of $3$.
\end{itemize}

A rigorous proof would require:
\begin{enumerate}
  \item An explicit $3$-adic formula for $L(\psi^4, 2)$ via Katz's $p$-adic
    CM $L$-functions \cite[Chapter III]{CITE-NEEDED-Katz1978},
    extending the framework of~\cite{Kriz2018};
  \item A Galois-cohomological computation controlling the $3$-adic valuation
    of the Euler factor at inert primes for symmetric power lifts of Hecke
    characters, along the lines of~\cite{BuyukbodukLei2016}.
\end{enumerate}
Such a rigorous proof is in preparation.

%% ============================================================
\section{Conjecture M184: twin-pair $c^4 = 4$ structure}
\label{sec:M184}
%% ============================================================

\subsection{Statement}

\begin{conjecture}[M184 --- twin-pair structure]
  \label{conj:M184}
  Let $K = \Q(\sqrt{D})$ be an imaginary quadratic field of class number $1$.
  Let $\sigma, \bar\sigma : K \hookrightarrow \C$ be the two complex embeddings.
  The corresponding Hecke characters $\psi_\sigma, \psi_{\bar\sigma}$ of
  infinity type $(4,0)$ give rise (in general) to two periods
  $\Omega_\sigma, \Omega_{\bar\sigma}$.
  Write $\Omega_{\bar\sigma} = c \cdot \Omega_\sigma$ for an algebraic $c$.
  Then:
  \begin{enumerate}
    \item If $D \in \{-11, -19\}$, then $c^4 = 4$ (i.e.\ $c = \sqrt{2}$ up to a
      root of unity), and the corresponding $\alpha_2$ values form a twin pair
      with ratio $\alpha_2'/\alpha_2 = c^4 = 4$.
    \item For all other Heegner discriminants
      $D \in \{-4, -3, -7, -43, -67, -163\}$, we have $c^4 = 1$
      (i.e.\ $\Omega_{\bar\sigma} = \pm\Omega_\sigma$), and $\alpha_2$ is unique.
  \end{enumerate}
\end{conjecture}

\begin{table}[h]
  \caption{Numerical verification of Conjecture~\ref{conj:M184}}
  \label{tab:M184}
  \renewcommand{\arraystretch}{1.25}
  \begin{tabular}{lrlrll}
    \toprule
    $K$ & $D$ & $\alpha_2$ pair & ratio & $c^4$ & Twin predicted? \\
    \midrule
    $\Q(i)$          & $-4$   & $\{1/12\}$         & ---   & $1$   & no (single) \\
    $\Q(\sqrt{-3})$  & $-3$   & $\{3/8\}$          & ---   & $1$   & no (single) \\
    $\Q(\sqrt{-7})$  & $-7$   & $\{28/3\}$         & ---   & $1$   & no (single) \\
    $\Q(\sqrt{-11})$ & $-11$  & $\{9/11,\,36/11\}$ & $4$   & $4$   & YES \\
    $\Q(\sqrt{-19})$ & $-19$  & $\{13/57,\,52/57\}$& $4$   & $4$   & YES \\
    $\Q(\sqrt{-43})$ & $-43$  & $\{214/129\}$      & ---   & $1$   & no (single) \\
    $\Q(\sqrt{-67})$ & $-67$  & $\{1519/201\}$     & ---   & $1$   & no (single) \\
    $\Q(\sqrt{-163})$& $-163$ & $\{196216792/3\}$  & ---   & $1$   & no (single) \\
    \bottomrule
  \end{tabular}
\end{table}

Table~\ref{tab:M184} shows agreement for all $8$ Heegner discriminants
($10$ individual $\alpha_2$ values), with twin pairs appearing exactly for
$D \in \{-11, -19\}$ as predicted ($2/2$ pairs verified).

\subsection{Class-field-theoretic interpretation}

The proposed equivalence~\emph{$c^4 = 1$ if and only if $E_D/\Q$
admits a rational $2$-isogeny} requires a careful cross-check against
LMFDB isogeny-class data. In particular, for $D = -11$ (LMFDB class
\texttt{121.b}) and $D = -19$ (LMFDB class \texttt{361.a}) the
elliptic curve with CM by $\OK_D$ does belong to a non-trivial
isogeny class containing a rational $2$-isogeny, so a direct
$c^4=4 \Leftrightarrow$ ``has rational $2$-isogeny'' equivalence is
consistent with these two anchors. We retain Conjecture~\ref{conj:M184}
as a numerical statement at $c^4$; a uniform class-field-theoretic
mechanism via the relative Hecke character $\psi_{\bar\sigma}/\psi_\sigma$
remains open.

Nevertheless, the numerical twin‑pair pattern for $D \in \{-11, -19\}$
remains unchanged and robust (Table~\ref{tab:M184}).
Consequently Conjecture~\ref{conj:M184} is \emph{downgraded} to a numerically
observed phenomenon whose underlying arithmetic mechanism is currently unknown.
Two candidate hypotheses are under investigation:
\begin{enumerate}
  \item The $2$-adic valuation of the period‑quadratic‑form discriminant
        \[
          \mathrm{disc}(Q) = \operatorname{Nm}_{K/\Q}(\Omega_{\bar\sigma}/\Omega_\sigma)^4,
        \]
        which can take the value $2$ (hence $c^4 = 4$) when the
        different of $K$ has a certain behaviour above~$2$.
  \item The conductor of the relative Hecke character
        $\psi_{\bar\sigma}/\psi_\sigma$ at primes above~$2$; if the
        conductor is non‑trivial, it forces an extra factor in the
        period ratio.
\end{enumerate}
A satisfactory answer is likely to emerge from a detailed study of the
mod‑$2$ Galois representation attached to $\psi$, but this requires new
computational data and is left for future work.

%% ============================================================
\section{The $\alpha_2$ invariant as a norm of Siegel elliptic units}
\label{sec:siegel}
%% ============================================================

The rationality and explicit form of $\alpha_2(D)$ admit a structural
interpretation combining the classical results of Damerell~\cite{Damerell1971},
Yager~\cite{Yager1982}, and Goldstein--Schappacher~\cite{GoldsteinSchappacher1981}
with the modern machinery of Kriz~\cite{Kriz2018}.

\subsection{The algebraic part as a Siegel unit norm}

Let $\psi_4$ be the Hecke character of $K$ of infinity type $(4,0)$ and
conductor $\mathfrak{f}$, giving rise to the CM newform $f_D$.
Let $K(\mathfrak{f})$ denote the associated ray class field.
For class-number-$1$ fields with trivial conductor ($\mathfrak{f} = (1)$,
equivalently $N = |D|$), the relevant ray class field is the field of
$N$-torsion points $K(N)$.

By combining the Kronecker limit formula with the Damerell--Yager explicit
evaluation, one obtains an identity of the form
\begin{equation}
  \label{eq:siegel}
  \frac{L(\psi_4, 2)}{\Omega_D^4} \;=\; \sum_{a \in \Cl(K, \mathfrak{f})}
    r_a \log |g_a|_v \pmod{\text{torsion}},
\end{equation}
where $g_a \in K(\mathfrak{f})^{\times}$ are Siegel elliptic units indexed by
ideal classes, and $r_a \in \Q$ are rational exponents determined by $\psi_4$.
Exponentiating~\eqref{eq:siegel}, the algebraic part $\alpha_2(D)$ is
(up to a root of unity) the norm from $K(\mathfrak{f})$ to $\Q$ of a product
of these units raised to rational powers:
\begin{equation}
  \label{eq:siegelnorm}
  \alpha_2(D) \;=\; \frac{1}{w} \cdot
    \Nm_{K(\mathfrak{f})/\Q}\!\left(\prod_{a \in \Cl(K,\mathfrak{f})}
      g_a^{r_a}\right),
\end{equation}
where $w = |\OK^{\times}|$ is the number of roots of unity in $\OK$.
Formula~\eqref{eq:siegelnorm} is a direct consequence of the formalism
developed in~\cite{GoldsteinSchappacher1981}; see also~\cite{Kriz2018}
for the $p$-adic analogue; cf.\ \cite{CITE-NEEDED-Katz1978}.

\subsection{Numerator arithmetic for $D = -163$}

The largest Heegner discriminant $D = -163$ provides the most striking
illustration.
One has $\alpha_2(-163) = 196\,216\,792 / 3$, with numerator factoring as
\[
  196\,216\,792 = 2^3 \cdot 163 \cdot 150\,473,
\]
where $150\,473$ is prime.
The factor $|D| = 163$ in the numerator aligns with the prediction from
formula~\eqref{eq:siegelnorm}: the norm from $K(163)$ to $\Q$ of the Siegel
units introduces a factor of $163$ via the Euler factor at the prime $163$
in the $L$-function of the CM elliptic curve.
The factor $8 = 2^3$ arises from the Bernoulli--Hurwitz number denominators
associated to the CM type of $\OK_{-163}$.
The prime $150\,473$ requires further arithmetic investigation to relate to
the class-field structure.

%% ============================================================
\section{Verification methodology}
\label{sec:verify}
%% ============================================================

\subsection{Software and precision}

All computations were performed with PARI/GP 2.15.4 \cite{PARI215}
using the \texttt{lfun} module for analytic continuation of $L$-functions
and the \texttt{bestappr} function for rational recognition.
Verification was carried out on a Vast.AI instance (AMD EPYC 80-core, $96$~GB
RAM) running PARI scripts in parallel across all Heegner discriminants.

\subsection{Critical lesson: partial sums fail for weight 5}

A key methodological point deserves emphasis.
For a weight-$5$ $L$-function $L(f,s) = \sum a_n n^{-s}$, the Dirichlet series
converges absolutely only for $\mathrm{Re}(s) > 3$.
At $s = 2$, which lies strictly inside the critical strip $[1,4]$, partial
sums $\sum_{n \leq X} a_n n^{-2}$ oscillate and do \emph{not} converge to
$L(f,2)$.
Any numerical claim about $\alpha_2(D)$ based on partial sums is therefore
unreliable.
All values in Theorem~\ref{thm:hierarchy} were computed via PARI's
\texttt{lfun} package, which implements analytic continuation via the
functional equation and integral representations.

\subsection{Negative results}

\paragraph{Class-number-$3$ fields.}
A sweep over $D \in \{-23, -31, -59, -83, -107, -139, -211, -283\}$
(the eight imaginary quadratic fields of class number $3$)
found \emph{zero} rational CM weight-$5$ newforms at any level up to $400$.
This strongly suggests the hierarchy of Theorem~\ref{thm:hierarchy} is a
genuinely class-number-$1$ phenomenon.

\paragraph{Higher weights.}
A parallel sweep over weights $k \in \{3, 5, 7, 9, 11, 13, 15, 17\}$
identified $18+$ additional rational CM newforms, concentrated at
weight $5$ and weight $9$.
Their algebraic parts will be studied in a sequel.

\subsection{Reproducible PARI script}

A fully reproducible script for computing $\alpha_2(D)$ at any imaginary
quadratic field is provided in Appendix~\ref{sec:pari}.

%% ============================================================
\section{Open problems}
\label{sec:open}
%% ============================================================

\begin{enumerate}
  \item \textbf{Rigorous proofs of M183 and M184.}
    Both conjectures are numerically verified but not proved.
    The denominator split conjecture (M183) should follow from a Galois-cohomological
    computation of the $3$-adic valuation of the Euler factor of
    $L(\mathrm{Sym}^4\psi, 2)$ at inert primes, using the $p$-adic interpolation
    machinery of~\cite{Kriz2018} and the Iwasawa-theoretic
    framework of~\cite{BuyukbodukLei2016}.
    The twin-pair conjecture (M184) is now recognized as a purely numerical
    pattern; the original expectation that it follows from the rational
    $2$-isogeny structure of $E_D$ has been falsified by a Cremona/LMFDB
    cross‑check (see internal note \texttt{HA\_B\_R2\_falsifies\_M184.md}),
    and the arithmetic mechanism behind the twin pairs remains unknown.

  \item \textbf{Extension to $h(K) > 1$.}
    Do rational CM weight-$5$ newforms exist for class-number-$> 1$ imaginary
    quadratic fields?
    Our computational evidence (zero forms found for all $h=3$ discriminants)
    suggests the answer is no, but a theoretical explanation is lacking.
    A necessary condition for rationality is that the Hecke character $\psi$
    be self-conjugate up to a finite-order idele-class character, which imposes
    strong restrictions beyond the class-number-$1$ case.

  \item \textbf{Higher-weight hierarchy.}
    The analogue of $\alpha_k(D) = L(f_D, \lfloor k/2 \rfloor) \cdot
    \pi^{\lfloor k/2 \rfloor} / \Omega_D^{k-1}$ for weight $k \in \{7, 9, 11,
    13, 15, 17\}$ at the same Heegner discriminants deserves systematic study.
    A Damerell--Yager-type formula for these values would constitute a complete
    arithmetic table of CM $L$-values over $h=1$ imaginary quadratic fields.

  \item \textbf{Formalisation.}
    The identity $\alpha_2(-4) = 1/12$ (Theorem~\ref{thm:hierarchy}, first row)
    has been partially axiomatised in Lean 4 (Appendix~\ref{sec:lean});
    a complete machine-checked proof requires Maass--Shimura operator theory and
    Chowla--Selberg formula infrastructure not yet available in Mathlib~4.

  \item \textbf{The numerator prime $150\,473$ for $D = -163$.}
    The prime factor $150\,473$ of the numerator $196\,216\,792$ does not
    appear in obvious class-number-formula or ring-class-field computations
    for $K(-163)$; its arithmetic origin is unclear and warrants investigation.
\end{enumerate}

%% ============================================================
\section*{Acknowledgments}
%% ============================================================

The author thanks the LMFDB team for the open-access modular forms database.
Computations were performed on Vast.AI cloud infrastructure.
Hypothesis generation and preliminary analysis were assisted by multi-LLM
tooling (LLM provider LLM and LLM tooling); all mathematical claims
were independently verified via PARI/GP and LMFDB lookup before inclusion.

%% ============================================================
%% BIBLIOGRAPHY
%% ============================================================

\begin{thebibliography}{99}

\bibitem{Damerell1971}
  R.~M. Damerell,
  \textit{$L$-functions of elliptic curves with complex multiplication, I},
  Acta Arith.\ \textbf{17} (1970), 287--301; \textit{II}, Acta Arith.\ \textbf{19}
  (1971), 311--317.

\bibitem{Yager1982}
  R.~I. Yager,
  \textit{On the special values of Hecke $L$-functions},
  Ann.\ of Math.\ \textbf{115} (1982), 219--265.

\bibitem{GoldsteinSchappacher1981}
  C.~Goldstein and N.~Schappacher,
  \textit{Séries d'Eisenstein et fonctions $L$ de courbes elliptiques à
  multiplication complexe},
  J.\ Reine Angew.\ Math.\ \textbf{327} (1981), 184--218.

\bibitem{Kriz2018}
  D.~Kriz,
  \textit{A new $p$-adic Maass--Shimura operator and supersingular
  Rankin--Selberg $p$-adic $L$-functions},
  arXiv:1805.03605, 2018.

\bibitem{BuyukbodukLei2016}
  K.~Büyükboduk and A.~Lei,
  \textit{Anticyclotomic $p$-ordinary Iwasawa theory of elliptic modular forms},
  arXiv:1602.07508, 2016.

\bibitem{KannoWatari2020}
  K.~Kanno and T.~Watari,
  \textit{$W=0$ complex structure moduli stabilization on CM-type
  K3 $\times$ K3 orbifolds},
  arXiv:2012.01111, 2020.

\bibitem{LMFDB2026}
  The LMFDB Collaboration,
  \textit{The $L$-functions and Modular Forms Database},
  \url{https://www.lmfdb.org}, accessed 2026.

\bibitem{PARI215}
  The PARI Group,
  \textit{PARI/GP, version 2.15.4},
  Univ.\ Bordeaux, 2024,
  \url{https://pari.math.u-bordeaux.fr/}.

%% References below are CITE-NEEDED: arXiv IDs proposed by LLM tooling
%% were verified as fabrications (wrong papers). Journal refs are correct
%% but no preprint IDs have been independently verified.

\bibitem{CITE-NEEDED-Katz1978}
  N.~M. Katz,
  \textit{$p$-adic $L$-functions for CM fields},
  Invent.\ Math.\ \textbf{49} (1978), 199--297.
  % CITE-NEEDED: no arXiv; journal citation unverified; DS-proposed arXiv:0704.2100
  % confirmed WRONG (galaxy photometry paper). Use with caution.

\bibitem{CITE-NEEDED-Kaneko2004}
  M.~Kaneko,
  \textit{A generalization of the Chowla--Selberg formula
  and special values of L-functions},
  J.\ Math.\ Soc.\ Japan (precise volume to be verified).
  % CITE-NEEDED: journal/volume requires independent verification;
  % previously suggested arXiv ID was wrong.

\bibitem{CITE-NEEDED-Yang2003}
  T.~Yang,
  \textit{CM values of modular functions and special $L$-values}.
  % CITE-NEEDED: DS-proposed arXiv:math/0305047 confirmed WRONG
  % (set theory paper by Schindler). Full citation requires independent lookup.

\bibitem{CITE-NEEDED-Brunault2016}
  F.~Brunault,
  \textit{Regulators for Rankin--Selberg products of modular forms}.
  % CITE-NEEDED: DS-proposed arXiv:1604.02683 confirmed WRONG
  % (telecommunications paper). Correct arXiv ID requires independent lookup.

\end{thebibliography}

%% ============================================================
\appendix
%% ============================================================

\section{PARI/GP script for $\alpha_2$ computation}
\label{sec:pari}

The following PARI/GP~2.15.4 script computes $\alpha_2(D)$ at any imaginary
quadratic field $K = \Q(\sqrt{D})$ of class number~$1$, using the
\texttt{lfun} module for analytic continuation and \texttt{bestappr} for
rational recognition.
Precision is set to $150$ decimal digits; increase \texttt{prec} for larger
discriminants.

\begin{verbatim}
/* alpha2_cm.gp  --  PARI/GP 2.15.4                                */
/* Compute alpha_2 = L(f,2)*Pi^2/Omega^4 for CM weight-5 newforms  */

alpha2_cm(D, N, prec=150) = {
  my(H, ff, f, L, val, omega, alpha, rat);
  localprec(prec);

  \\ 1. Find weight-5 newform of level N with CM by Q(sqrt(D))
  H  = mfinit([N, 5, Mod(kronecker(D,x), x^2-1)], 1);
  ff = mfeigenbasis(H);
  f  = [];
  for (i = 1, #ff,
    if (mfisCM(ff[i]) && mfcmfield(ff[i]) == D,
      f = ff[i]; break));
  if (#f == 0, error("No CM form found for D=", D, " N=", N));

  \\ 2. L-value L(f,2) via analytic continuation
  \\ Note: lfunmf() is available in PARI >= 2.15 (replaces lfuninit+mf).
  \\ If unavailable, construct via lfuncreate on the Hecke character directly.
  L   = lfunmf(f);
  val = lfun(L, 2);

  \\ 3. Chowla-Selberg period Omega_D
  omega = chowla_selberg(D);

  \\ 4. alpha_2 = L(f,2)*Pi^2/Omega^4
  alpha = val * Pi^2 / omega^4;

  \\ 5. Rational recognition
  rat = bestappr(alpha, 10^(prec - 10));
  printf("D=%d  alpha_2 = %s  denom=%d\n", D, rat, denominator(rat));
  return(rat);
}

chowla_selberg(D) = {
  my(absD=abs(D), chi=kronecker(D,), w, prod=1.0);
  w = if(D==-3, 6, D==-4, 4, 2);
  forstep(a=1, absD-1, 1,
    if(gcd(a, absD)==1,
      prod *= gamma(a/absD)^(w/4 * chi(a))));
  return(prod / sqrt(2*Pi));
}

\\ Example calls (matches Theorem 2.1):
\p 150
alpha2_cm(-4,  4)   \\ expects 1/12
alpha2_cm(-3,  12)  \\ expects 3/8
alpha2_cm(-7,  7)   \\ expects 28/3
alpha2_cm(-11, 11)  \\ expects 9/11 (first curve)
alpha2_cm(-19, 19)  \\ expects 13/57 (first curve)
alpha2_cm(-43, 43)  \\ expects 214/129
alpha2_cm(-67, 67)  \\ expects 1519/201
alpha2_cm(-163,163) \\ expects 196216792/3
\end{verbatim}

\begin{remark}
  The \texttt{mfisCM} and \texttt{mfcmfield} PARI functions require
  PARI/GP $\geq 2.15$.
  For $D = -11$ and $D = -19$, two distinct elliptic curves with CM by
  $\OK_D$ exist over $\Q$; run the script with each curve separately
  (specifying the Weierstrass coefficients via \texttt{ellinit}) to recover
  both values of the twin pair.
\end{remark}

\section{Lean 4 formalisation stub}
\label{sec:lean}

A partial formalisation of the $\alpha_2(-4) = 1/12$ identity
(Theorem~\ref{thm:hierarchy}, row~1) in Lean~4 / Mathlib~4 has been
written as a stub with all theorems closed by \texttt{sorry}.
The stub is available at:
\begin{center}
  \texttt{/root/crossed-cosmos/lean/EciM142.lean}
\end{center}
(or, upon publication, at the paper's GitHub repository).

The stub uses real Mathlib~4 API where available
(\texttt{WeierstrassCurve}, \texttt{EllipticCurve}, \texttt{ModularForm},
\texttt{LSeries}, \texttt{CongruenceSubgroup.Gamma0}, \texttt{Real.pi},
\texttt{Complex.Gamma}) and axiomatises the objects for which Mathlib~4
infrastructure is not yet available (the CM period $\Omega$, the
Maass--Shimura operator, the functional equation).
A complete machine-checked proof would require:
\begin{enumerate}
  \item Analytic continuation of $L(f,2)$ in Mathlib (not yet available);
  \item The Chowla--Selberg formula for the CM period (not yet available);
  \item The explicit rational factor from the CM type of $\Z[i]$
    and the evaluation $\Gamma(1/4)^2 = 4\sqrt{2\pi}\,\Omega$.
\end{enumerate}
We declare this as a long-term formalisation target for the Mathlib community.

\end{document}
